\section{作品概述}
\subsection{项目背景}
随着信息技术的发展，数据安全不仅关系到个人隐私和企业商业秘密，也是国家安全的重要组成部分。
传统数据加密技术已广泛应用于信息传输领域，能够有效保护传输内容，但难以隐藏通信行为本身。在实际网络环境中，
即使通信内容经过加密，第三方仍然可能通过分析通信时间、流量大小等推测用户通信关系以及活动规律，从而造成隐私泄露。
因此，数据传输不仅需要保障信息内容的安全性，还需要关注敏感信息传输的隐蔽性。

为满足这一要求，图像隐写术通过将秘密信息隐藏在自然图像中，降低敏感信息传输被发现的可能性\supercite{1-wang-moulin-2008,2-yang-et-al-2024}。
然而，部分基于深度学习的图像隐写方案需要额外训练专用网络，增加了模型准备和部署的负担。例如，例如，HiDDeN通过联合训练编码网络和解码网络，实现秘密信息的嵌入与恢复\supercite{3-zhu-et-al-2018}。
SteganoGAN则采用生成对抗网络，通过联合优化编码器、解码器和判别网络，提高图像隐写的信息容量和视觉质量\supercite{4-zhang-et-al-2019-steganogan}。
这些方案虽然在信息隐藏方面取得了很好效果，但其需要专用网络的训练与部署工作。因此，如何充分复用现有资源，在减少额外模型准备成本的同时实现安全且隐蔽的信息传输，是一个关键问题。

深度神经网络（Deep Neural Network，DNN）已广泛应用于图像处理等领域，并逐渐成为各类信息系统的重要计算资源\supercite{5-lecun-et-al-2015}。
近年来，DNN与密码学\supercite{6-abadi-andersen-2016,7-gerault-et-al-2026},图像隐写\supercite{8-luo-et-al-2023}等领域的交叉研究，也表明
DNN不仅能够用于传统的图像处理任务，还可以在信息安全领域发挥重要作用。从实现原理来看，视觉DNN能够通过非线性变换，将输入图像映射为高维神经元状态，适合表示含义丰富的信息\supercite{9-bengio-et-al-2013}。
同时，DNN的神经元状态会随着输入图像的变化而变化，因此可以通过调整输入图像来控制神经元状态。我们受此启发，提出了一种基于表示的框架。该框架可以不修改现有DNN参数和原有功能，利用其内部神经元状态实现
安全且隐蔽的信息表示。具体而言，我们将信息标识为选定神经元的目标状态，并通过对抗优化将随机图像迭代为外观自然的密文图像。当密文图像与预共享图像密钥结合时，可在DNN中诱导出目标状态，从而实现密钥控制的
消息恢复。

然而，此类框架仍面临两个主要挑战，首先，如何准确稳定地诱导DNN产生目标神经元状态。现有视觉DNN主要面向图像分类等任务，其内部神经元
由多层网络共同决定，难以直接控制。已有研究表明，通过对术后如图像施加细微的对抗扰动，可以改变DNN的输出结果\supercite{10-szegedy-et-al-2014,11-goodfellow-et-al-2015}。
但是由于不同神经元对输入扰动的敏感程度存在差异\supercite{12-zhang-et-al-2019-sensitivity,13-suri-evans-2020}，部分神经元的相应可能随着输入变化发生明显干瘪，而另一些神经元的相应变化小，
难以达到预设的目标值，我们称其为不响应神经元。当利用神经元表示消息时，不响应神经元可能导致部分消息比特无法正确编码，从而降低信息表示的准确性，并增加生成过程所需的迭代次数。
其次，如何在诱导目标神经元状态的同时抑制密文图像的信息泄露。为了诱导DNN神经元至目标状态，需要通过对抗优化不断调整输入图像。但是图像修改可能导致密文图像与原始图像存在视觉差异，攻击者有可能通过
图像分析发现异常，即直接视觉泄露\supercite{14-hu-et-al-2024}。并且由于DNN的结构和参数公开，攻击者可以将截获密文直接输入相同DNN，分析其内部神经元状态异常变化，形成间接DNN状态泄露隐私泄露。
总而言之，提高目标状态诱导的准确性与稳定性，消除整体视觉特征和内部DNN状态异常，是设计方案的两大主要困难。

\subsection{相关工作}
在神经网络与密码学的交叉研究中，已有工作主要探索利用神经网络学习信息保护机制或实现密码操作。
Abadi和Andersen\supercite{6-abadi-andersen-2016}提出了对抗神经密码学，通过联合训练负责加密、解密和攻击的三个神经网络，使网络自主学习信息保护方法，但该过程需要专门训练模型。
Gérault等人\supercite{7-gerault-et-al-2026}研究了如何利用DNN实现密码操作，针对传统密码算法的离散计算与神经网络连续计算之间的差异，提出了基于ReLU的实现方法，并分析了其正确性与安全性。这些研究表明了神经网络在密码学领域的应用潜力，但与直接利用现有DNN进行隐蔽信息表示仍有所不同。

在图像隐写领域，已有研究主要关注信息隐藏容量、图像质量和鲁棒性等问题\supercite{2-yang-et-al-2024,16-yao-et-al-2024,17-li-et-al-2025,18-liu-et-al-2026}。
HiDDeN\supercite{3-zhu-et-al-2018}通过联合训练编码网络和解码网络实现信息嵌入与恢复，SteganoGAN\supercite{4-zhang-et-al-2019-steganogan}则利用生成对抗网络提高信息隐藏容量，但这些方法需要额外训练专用网络，增加了模型准备和部署成本。
为减少对专用网络训练的依赖，Kishore等人\supercite{19-kishore-et-al-2021}提出固定神经网络隐写，通过优化输入图像而不修改网络参数实现信息隐藏。Ghamizi等人\supercite{20-ghamizi-et-al-2021}利用对抗扰动改变神经网络的分类结果，将秘密信息表示为目标类别。
Luo等人\supercite{8-luo-et-al-2023}进一步在固定神经网络隐写中引入密钥控制，使消息恢复受到密钥约束，并优化隐写图像的视觉质量。
\subsection{作品特色}
通过分析现有方案的不足，我们设计了一种基于图像密钥控制DNN状态的安全隐蔽信息表示框架（Secure and Stealthy Information Representation，SSIR）。该方案不需要重新训练神经网络，基于图像密钥融合、分组编码和对抗优化实现，具有以下特点。

(1)新型信息表示方法。与直接将信息嵌入图像的方式不同，SSIR将消息编码为DNN中选定神经元的目标状态。发送方通过对抗优化调整图像，使其与图像密钥结合后能够产生对应的神经元状态，
接收方使用相同密钥和DNN恢复消息。同时，考虑到部分神经元对输入扰动不敏感，我们采用分组编码，对相邻神经元的响应取平滑平均，减少单个神经元不响应带来的影响。

(2)可靠低泄露。在生成密文图像时，对抗扰动可能留下视觉痕迹，也可能使DNN的内部状态发生异常变化。
针对这两种情况，我们分别设计了直接视觉泄漏损失和间接DNN状态泄漏损失，并将其加入图像优化过程，尽量减少攻击者在没有图像密钥时能够观察到的信息痕迹。 我们在ImageNet数据集和真实图像上进行了测试。
实验中，通过调整学习率，直接视觉泄漏指标最高降低99.79\%；
加入间接泄漏损失后，神经元状态偏差平均降低56.02\%。对于150位消息，生成密文图像约需300～500 ms，恢复消息约需10 ms。

（3）查询有界安全性分析。 除实验测试外，我们还对方案进行了理论分析，包括图像密钥与DNN非线性响应的关系、分组编码后的恢复一致性、密文图像随机生成中的重复输出可能，以及同一图像在兼容密钥下的多语义解释。
在此基础上，推导了查询有界的IND-CPA和IND-CCA1区分优势上界。
其中，查询有界指攻击者请求生成或恢复的次数有限，IND-CPA对应攻击者不能请求恢复时仍难以区分挑战密文，IND-CCA1则进一步计入挑战前恢复查询带来的额外区分差距。
分析表明，在满足相应分布条件并限制每密钥查询次数时，攻击者的区分优势不超过给定上界。该结论明确了安全成立的条件，并为实际部署中确定每密钥查询上限、实施访问限制和制定密钥更新策略提供依据。

（4）模型兼容与参数指导。 SSIR无需修改DNN参数，也不需要额外训练专用模型。它可以直接复用现有DNN完成信息表示于回复，降低在既有DNN中集成的门槛。本文跨设备测试表明，只要生成段和恢复段统一数值处理规则，即可无误收发消息，为实际跨设备
部署提供明确工程指导。分组编码用于缓解部分神经元造成的生成失败，实验表明分组2再可表示消息量与生成速度之间较为均衡。本文形成了可直接用于部署的分组大小、学习率和跨设备数值规则选择参考，使方案
能够在不改动原有模型的前提下降低部署与调参成本，并按应用需求在成功率、图像质量和生成时延之间进行权衡。

\subsection{应用前景分析}
当前，敏感信息传输面临双重困境：传统加密技术虽然能够保护消息内容，但无法隐藏通信行为本身，
密文在公开信道上的频繁出现容易引起第三方注意；图像隐写技术虽然能够把消息隐藏到自然载体中，却往往依赖专用网络训练或额外组件，
部署成本较高。随着即时通信、物联网等场景的发展，会话密钥、控制指令等短小敏感信息需要在公开信道上
频繁、快速、隐蔽地传递。SSIR 正是面向这一需求提出的，它利用图像密钥控制现有公开 DNN 的神经元状态，
将消息编码为目标状态并生成外观自然的密文图像，接收方叠加密钥并输入 DNN 即可恢复。在具体实际部署中，SSIR 具有以下典型应用价值：

在即时通信场景中，传统的端到端加密虽然保护了聊天内容，但在建立会话时，密钥协商的报文特征往往非常明显，容易被流量审计设备识别并标记。
SSIR可以使得双方直接复用手机或个人电脑上现有的轻量级视觉模型，例如手机相册自带的分类模型，将新的会话密钥编码进一张普通的日常照片中，
发送方只需像发送普通图片一样将密文图像发给对方，接收方叠加与共享的图像密钥即可瞬间恢复出新密钥。这样将密钥分发隐蔽成图一篇分享，
大大降低密钥传递过程被检测的风险，且无需再终端额外部署庞大的加密网关或训练专用隐写模型。

在物联网场景中，设备间经常需要下发开启阀门、关闭传感器等短小控制指令。如果指令以明文或规定格式的密文传输，极易被监听者识别规律，甚至遭到
针对性重放或阻断。由于边缘设备算力有限，无法运行复杂的隐写网络，SSIR可以复用设备端已有的、用于图像质检或环境感知的DNN模型。下发端将指令
编码为外观自然的密文图像，接收端直接从DNN神经元状态中读取指令。即使信道被持续监听，攻击者看到的也只是一组带有轻微自然噪点的图像流，
难以区分哪一张包含了关键动作指令，从而提升了控制信道的抗干扰与抗流量分析能力。

在实际企业运行时，部署方可以根据自身对安全强度的要求，确定每个图像密钥允许的生成与恢复次数上限。系统后台对同一密钥的会话进行累计计数，
当达到安全阈值前，自动触发密钥更新，各方独立生成并预共享新的替换密钥。系统会同时限制外部攻击者获取无限次的恢复响应。这一机制符合既定的
密钥管理指南，避免了因单密钥过度使用而带来的安全风险，使系统在长期运行中保持了安全底线。

综上所述，SSIR不仅是一种底层的信息表示方法，更具备显著的工程落地价值。它既能作为现有加密通信的增强曾，服务于企业内部的密钥管理或客户通信，也能在
算力受限的物联网中，无需修改模型和额外训练，以极低的部署成本实现高可靠、低泄露的敏感信息传递。